\answer{ex:19:1}{(a)~シナプス1個で$6.7\mV$（$R=66.7$~M$\Omega$）なので、そもそも閾値に達するには$4$個以上が必要（3個では漸近的にしか届かない）。$10\ms$では$8$個、$5\ms$では$14$個。(b)~$0.1\ms\ll\tau_m$、漏れによる補正は$\sim0.25\%$。}
\answer{ex:19:2}{(a)~$1.75\cdot10^{10}$（ミキサの4分の1が何にでも応答する）、$4.4\cdot10^9$、$0.06$（$k=40$ではほとんど応答しない）。(b)~$10^3$倍：$2\cdot10^5$対$200$。}
\answer{ex:19:3}{(a)~二項係数の対称性により$C(2n,n)=2^{2n-1}$。(b)~$0.93$と$0.09$。$P$についての遷移幅は$\sim\sqrt{2n}$、相対幅は$\sim1/\sqrt n$。}
\answer{ex:19:4}{1本の樹状突起からでは、入力がいくつあっても$V_{\text{soma}}<\kappa E_E<\theta$。$g_0=0.5g_L$では各樹状突起に$n\ge3$個が必要。}
\answer{ex:19:5}{$I\ge C\sigma_V/\sigma_t=15$~nA、すなわち$150$個の同期シナプスが必要で、非現実的である。小さなニューロンなら$1$~nAで足り、$4$~nAのシナプスではジッタは$5$~\textmu s。}
\answer{ex:19:6}{近端での最大：$0.03\,\tau_m$で$0.97$（$L=0.2$）、$0.32\,\tau_m$で$0.66$（$L=1$）、$0.79\,\tau_m$で$0.33$（$L=2$）。全容量による見積もり~\eqref{eq:dendriteload}は$L\ll1$でしか正しくない。}
\answer{ex:19:7}{自由課題。$m$を固定すると、応答時間は記憶する$P$に比例して増える。割合$m=\varphi n$を固定すると$n$に依存しなくなる。大きなニューロンの遅さは、大きさそのものではなく、多くの入力を足し合わせることの帰結である。}
